\documentclass[10pt,letterpaper]{article}

% ---------- Packages ----------
\usepackage[margin=0.68in]{geometry}
\usepackage{enumitem}
\usepackage{tabularx}
\usepackage{titlesec}
\usepackage[hidelinks]{hyperref}
\usepackage{xcolor}
\usepackage{microtype}
\usepackage{fontawesome5}
\usepackage[backend=biber,style=numeric,sorting=ydnt,maxbibnames=99]{biblatex}
\addbibresource{publications.bib}

% ---------- Typography ----------
\definecolor{accent}{HTML}{1F4E79}
\setlength{\parindent}{0pt}
\setlength{\parskip}{0pt}
\pagestyle{empty}
\renewcommand{\familydefault}{\sfdefault}

\titleformat{\section}
  {\large\bfseries\color{accent}}
  {}{0pt}{}
  [\vspace{-0.35em}\titlerule]
\titlespacing*{\section}{0pt}{0.9em}{0.45em}

\setlist[itemize]{leftmargin=1.15em,itemsep=0.18em,topsep=0.2em,parsep=0pt}

% ---------- Reusable commands ----------
\newcommand{\cvname}[1]{\begin{center}{\Huge\bfseries #1}\end{center}}
\newcommand{\cvcontact}[1]{\begin{center}\small #1\end{center}}

\newcommand{\cventry}[4]{%
  \begin{tabularx}{\textwidth}{@{}Xr@{}}
    \textbf{#1} & \textbf{#2} \\
    \textit{#3} & \textit{#4}
  \end{tabularx}\vspace{0.15em}
}

\newcommand{\cvsubentry}[2]{%
  \begin{tabularx}{\textwidth}{@{}Xr@{}}
    #1 & #2
  \end{tabularx}
}

\newcommand{\cvskill}[2]{\textbf{#1:} #2\\}

% ---------- Document ----------
\begin{document}

\cvname{Xiang He}
\cvcontact{
  Ph.D. Candidate in Geosciences, The University of Texas at Austin\\[0.15em]
  \href{mailto:hexiang@utexas.edu}{\faEnvelope\ hexiang@utexas.edu}
  \quad | \quad
  \href{https://hx38324.github.io/}{\faGlobe\ hx38324.github.io}
  \quad | \quad
  \href{https://github.com/hx38324}{\faGithub\ github.com/hx38324}
  % \quad | \quad \href{https://orcid.org/YOUR-ORCID}{ORCID}
}

\section{Research Profile}
Computational geodynamicist studying how mantle flow, subduction, lithospheric deformation, and mantle rheology shape observable surface motions. I build physics-based numerical models and compare model predictions with geological and geophysical constraints, with current emphasis on Western U.S. vertical motion and Cascadia.

\section{Education}
\cventry{The University of Texas at Austin}{Aug. 2024--Present}
{Ph.D. in Geosciences / Geophysics}{Austin, TX, USA}
\begin{itemize}
  \item Advisor: Prof. Thorsten W. Becker
  \item Research focus: computational geodynamics, mantle flow, subduction, dynamic topography, vertical crustal motion
\end{itemize}

\cventry{University of Science and Technology of China}{Sep. 2020--Jun. 2024}
{B.S. in Geophysics}{Hefei, Anhui, China}
\begin{itemize}
  \item Guo Moruo Scholarship; China National Scholarship
\end{itemize}

\section{Research Experience}
\cventry{The University of Texas at Austin}{Aug. 2024--Present}
{Graduate Researcher, Institute for Geophysics / Department of Earth and Planetary Sciences}{Austin, TX}
\begin{itemize}
  \item Develop global and regional geodynamic models to investigate mantle-flow contributions to long-wavelength surface deformation and vertical motion.
  \item Current project: construct Cascadia subduction and Western U.S. geodynamic models and evaluate predicted vertical motion against geophysical observations.
  \item Design numerical experiments to test sensitivity to slab geometry, mantle rheology, density structure, and boundary conditions.
  \item Run and analyze large-scale simulations using ASPECT/CitcomS and high-performance computing workflows; process model output with Python, GMT, and related tools.
\end{itemize}

\cventry{University of Illinois Urbana-Champaign}{May 2023--Oct. 2023}
{Visiting Researcher, Department of Earth Science and Environmental Change}{Urbana, IL}
\begin{itemize}
  \item Conducted 3-D regional geodynamic modeling of intraplate deformation in East Asia under the supervision of Prof. Lijun Liu.
\end{itemize}

\cventry{University of Science and Technology of China}{Jan. 2023--Aug. 2023}
{Undergraduate Researcher, School of Earth and Space Sciences}{Hefei, China}
\begin{itemize}
  \item Modeled planetary crust formation and thermal-mechanical evolution under the supervision of Prof. Wei Leng.
\end{itemize}

\cventry{University of Science and Technology of China}{Apr. 2022--Jun. 2024}
{Undergraduate Researcher, School of Earth and Space Sciences}{Hefei, China}
\begin{itemize}
  \item Investigated plume--subduction interactions and the influence of African mantle dynamics on Tethyan evolution under the supervision of Prof. Jinshui Huang.
\end{itemize}

\section{Publications}
% Keep publication metadata ONLY in publications.bib.
% Accepted/in-press papers can use keywords={accepted} in the BibTeX entry.
\nocite{*}
\printbibliography[heading=none]

\section{Selected Research Themes}
\begin{itemize}
  \item \textbf{Western U.S. vertical motion and Cascadia:} subduction-system dynamics, mantle flow, dynamic topography, and comparison with surface observations.
  \item \textbf{Mantle rheology and plate motion:} sensitivity of plate-scale deformation and surface observables to mantle viscosity structure and density heterogeneity.
  \item \textbf{Subduction and plume dynamics:} numerical models of single/double subduction, plume--plate interactions, and long-term tectonic evolution.
\end{itemize}

\section{Technical Skills}
\cvskill{Programming}{Python, C, Fortran, MATLAB, shell scripting}
\cvskill{Geodynamic / Geophysical Software}{ASPECT, CitcomS, GPlates, GMT / PyGMT, ParaView}
\cvskill{Scientific Computing}{HPC workflows, MPI-based simulations, Linux, Git, numerical modeling, inverse problems, data--model comparison}
\cvskill{Methods}{Finite-element geodynamics, sensitivity analysis, geophysical data processing, visualization}

\section{Awards \& Honors}
\cvsubentry{Graduate Fellowship, The University of Texas at Austin}{2024--2025}
\cvsubentry{Guo Moruo Scholarship, USTC}{2024}
\cvsubentry{China National Scholarship (top 0.2\% in China)}{2023}
\cvsubentry{Yang Yongman Scholarship, USTC}{2023}
\cvsubentry{USTC Outstanding Student Scholarship, Gold Medal (top 3\%)}{2022}
\cvsubentry{USTC Outstanding Student Scholarship, Silver Medal (top 10\%)}{2021}
\cvsubentry{Zhao Jiuzhang Earth Science Honored Class Scholarship}{2021, 2022}

\section{Teaching}
\cventry{University of Science and Technology of China}{Mar. 2023--Jun. 2023}
{Teaching Assistant, Electromagnetism A (PHYS1004A.04.2023SP)}{Hefei, China}
Instructor: Prof. Gan Qin

% ---------- Optional future sections ----------
% \section{Presentations}
% \section{Professional Service}
% \section{Open-Source Software / Data}
% \section{Professional Memberships}
% \section{References}

\end{document}
